\documentclass[exercice,a4]{thedocument}%
\usepackage{texmaths-tikz}%
\usepackage{pstricks}%
\usepackage{pst-bezier}%
\usepackage{pstricks-add}%
\usepackage{pst-eucl}%
\usepackage{pst-func}%
\usepackage{pst-tree}%
\usepackage{pst-circ}%
\usepackage{tikz}%
\usepackage{pgfplots}%
\usepackage{tkz-tab}%
\usepackage{tikzpagenodes}%
\usepackage{tabularray}%
\usepackage{tabularx}%
\usepackage{multirow}%
\usepackage{multicol}%
\usepackage{scratch3}%
\usepackage{enumitem}%
\usepackage[np]{numprint}%
\usepackage{pas-tableur}%
\usepackage{listings}%
\usepackage{graphicx}%
\usepackage{adjustbox}%
\usepackage{needspace}%
\usepackage{xcolor}%
\usepackage{eurosym}%
\usepackage{diagbox}%
\usepackage{colortbl}%
\usepackage{multido}%
\usepackage{icomma}%
\usepackage[french]{babel}%
\startdatakeys
    \source{26GENMATG11}
    \BO{2026}
    \cycle{4}
    \competencesDuSocle{CA,RA,CO}
    \objectifsApprentissage{C4A1aa,C4A2ad,C4B2cb,C4D2aa,C4D3ac}
    \motsCles{volume,cylindre,boule,conversion,proportionnalité}
\enddatakeys
\definecolor{gris}{gray}{0.5}%
\providecommand{\wideparen}[1]{\overset{\frown}{#1}}%
\begin{document}%
\providecommand{\degres}{^\circ}%
\emph{Les deux parties de l'exercice sont indépendantes.}

\subsection*{Partie A}

Une confiserie fabrique des bonbons multicolores au goût réglisse. Ces bonbons de longueur totale 15 mm ont la forme de gélules constituées de trois parties : un cylindre et deux demi-boules identiques de rayon $R = 2{,}5$ mm comme le montre le schéma ci-dessous.

\begin{minipage}{0.48\linewidth}
%%Début Tikz
\begin{center}
\begin{tikzpicture}[rotate=-35]
 \draw[line width=1.2pt] (-1.2,4.8) arc(180:0:1.2)--(1.2,0) arc (0:-180:1.2) -- cycle
 (-1.2,4.8) arc(-180:0:1.2 and 0.45)
 (-1.2,0) arc(-180:0:1.2 and 0.45);
 \draw[line width=1.2pt,dash pattern=on 6pt off 6pt]
 (-1.2,4.8) arc(180:0:1.2 and 0.45)
 (-1.2,0) arc(180:0:1.2 and 0.45);
 \draw[line width=0.8pt,dash pattern=on 3pt off 3pt]
 (-1.2,4.8) -- (0,4.8) node[pos=0.5,text width=1.2mm, align=center,rotate=-35,scale=0.8]{2,5 mm};
 \draw[<->, line width=0.8pt, > = stealth,shift={(-0.5,0)}] (-1.2,0)--(-1.2,4.8) node[pos=0.5,left]{$h$};
 \draw[<->, line width=0.8pt, > = stealth,shift={(0.5,0)}] (1.2,-1.2)--(1.2,6) node[pos=0.5,below, sloped,rotate=-35]{15 mm};
 \draw[line width=0.8pt] (-2pt,-2pt)--(2pt,2pt) (-2pt,2pt)--(2pt,-2pt);
 \draw[line width=0.8pt,shift={(0,4.8)}] (-2pt,-2pt)--(2pt,2pt) (-2pt,2pt)--(2pt,-2pt);
\end{tikzpicture}
\end{center}
%%Fin Tikz
%%%Début PSTricks
%\psset{unit=1cm,arrowsize=2pt 3}
%\begin{pspicture}(-2,-1.5)(2,6)
%	\psellipticarc(0,0)(1.1,0.4){180}{360}
%	\psellipticarc[linestyle=dashed](0,0)(1.1,0.4){0}{180}
%	\psarc(0,0){1.1}{180}{360}
%	\psline(1.1,0)(1.1,4.7)
%	\psline(-1.1,0)(-1.1,4.7)
%	\psellipticarc(0,4.7)(1.1,0.4){180}{360}
%	\psellipticarc[linestyle=dashed](0,4.7)(1.1,0.4){0}{180}
%	\psarc(0,4.7){1.1}{0}{180}
%	\psline[linewidth=0.8pt]{<->}(-1.5,0)(-1.5,4.7)\uput[l](-1.5,2.35){$h$}
%	\psline[linewidth=0.8pt]{<->}(1.5,-1.1)(1.5,5.8)\rput{90}(1.8,2.35){15 mm}
%	\psdots[dotstyle=+,dotangle=45](0,0)(0,4.7)
%	\psline[linestyle=dashed](0,4.7)(-1.1,4.7)\uput[u](-0.55,4.7){\footnotesize 2,5 mm}
%\end{pspicture}
%%%Fin PSTricks
%\emph{(Gélule : cylindre de hauteur $h$ flanqué de deux demi-sphères de rayon $2,5$ mm ; longueur totale $15$ mm.)}
\end{minipage}\hfill
\fbox{
	\begin{minipage}{0.4\textwidth}
	\textbf{Rappels}

	\begin{itemize}
		\item Volume d'un cylindre de rayon $R$ et de hauteur $h$ :
		\[ V = \pi \times R^2 \times h \]

		\item Volume d'une boule de rayon $R$ :
		\[ V = \frac{4}{3} \times \pi \times R^3 \]

		\item $1\text{ L} = 1\text{ dm}^3$
	\end{itemize}
\end{minipage}}\hfill~

\begin{enumerate}
	\item Montrer que la hauteur $h$ du cylindre est égale à $10$ mm.
	\item \begin{enumerate}
		\item Montrer que le volume de la partie cylindrique d'un bonbon est environ égal à $196$ mm$^3$.
		\item Léa affirme que le volume total d'un bonbon est compris entre $260$ et $262$ mm$^3$.

		A-t-elle raison ?
	\end{enumerate}
	\item Pour réaliser ces bonbons, la confiserie fabrique un mélange d'ingrédients qui est chauffé puis versé dans des moules en forme de gélule avant d'être refroidi.

	La confiserie fabrique chaque jour $83$L de mélange.

	Avec cette quantité de mélange, peut-elle produire plus de \np{300000} bonbons par jour ?
\end{enumerate}

\subsection*{Partie B}

Dans un magasin, les bonbons sont vendus en deux formats possibles :

\medskip
	\begin{tblr}{colspec={|X[c]|X[c]|},row{1}={font=\bfseries,cmd={\underline}}}
		\hline
		Format A          & Format B          \\
		Sachet de 500 g de bonbons & Sachet de 250 g de bonbons \\
		7,90~\euro{} le sachet     & 4,30~\euro{} le sachet     \\
		                           & \emph{\textbf{Offre promotionnelle} : pour 3 sachets achetés, le quatrième est à moitié prix.} \\ \hline
	\end{tblr}

\medskip
Léa veut acheter 1 kg de bonbons.

Quel format doit-elle choisir pour payer le moins cher possible ? \textbf{Argumenter la réponse en précisant la démarche.}

\needspace{10\baselineskip}
\medskip
\showthesource
\end{document}
